\documentclass[10pt, a4paper, twocolumn]{article}

% ── packages ────────────────────────────────────────────────────────────────
\usepackage[a4paper,
            top=22mm, bottom=24mm,
            left=15mm, right=15mm,
            columnsep=7mm]{geometry}
\usepackage[T1]{fontenc}
\usepackage[utf8]{inputenc}
\usepackage{microtype}
\usepackage{mathptmx}
\usepackage{parskip}
\usepackage{titlesec}
\usepackage{fancyhdr}
\usepackage{enumitem}
\usepackage{xcolor}
\usepackage{booktabs}
\usepackage{array}
\usepackage{hyperref}
\usepackage{setspace}
\usepackage{amssymb}
\usepackage{authblk}
\usepackage{abstract}
\usepackage{caption}

% ── colours ─────────────────────────────────────────────────────────────────
\definecolor{darktext}{HTML}{111111}
\definecolor{mutedgrey}{HTML}{555555}
\definecolor{rulegrey}{HTML}{222222}

% ── hyperref ────────────────────────────────────────────────────────────────
\hypersetup{
  colorlinks=false,
  hidelinks,
  pdftitle={Diotima: Compliance-Grade AI Infrastructure for Formative Assessment in Education},
  pdfauthor={Noval Consulting}
}

% ── typography ───────────────────────────────────────────────────────────────
\setlength{\parskip}{4pt}
\setlength{\parindent}{1em}
\frenchspacing

% ── section headings — centred small-caps, like Regional Studies ─────────────
\titleformat{\section}
  {\centering\normalsize\bfseries\scshape}
  {}
  {0pt}
  {}
  []

\titleformat{\subsection}
  {\normalsize\itshape}
  {}
  {0pt}
  {}

\titlespacing{\section}{0pt}{14pt}{6pt}
\titlespacing{\subsection}{0pt}{10pt}{3pt}

% ── abstract styling ─────────────────────────────────────────────────────────
\setlength{\absleftindent}{0pt}
\setlength{\absrightindent}{0pt}
\renewcommand{\abstractnamefont}{\normalsize\bfseries\scshape}
\renewcommand{\abstracttextfont}{\small}

% ── header / footer ──────────────────────────────────────────────────────────
\pagestyle{fancy}
\fancyhf{}
\renewcommand{\headrulewidth}{0pt}
\fancyhead[L]{\small\textit{AI \& Education Policy}}
\fancyhead[R]{\small\textit{Noval Consulting Working Paper}}
\fancyfoot[C]{\small\thepage}
\setlength{\headheight}{14pt}

% ── bullet list style ─────────────────────────────────────────────────────────
\setlist[itemize,1]{
  label={--},
  leftmargin=1.2em,
  itemsep=1pt,
  topsep=3pt,
  parsep=0pt
}

% ── title block — one column, then two-column body ───────────────────────────
\makeatletter
\renewcommand{\maketitle}{%
  \twocolumn[%
    \begin{@twocolumnfalse}
      % journal slug
      {\small\textit{AI \& Education Policy},
       Working Paper Series,\ pp.\ 1--12.}
      \vspace{6pt}
      \hrule height 0.8pt
      \vspace{2pt}
      \hrule height 0.2pt
      \vspace{14pt}

      % title
      {\centering
        {\fontsize{22}{27}\selectfont
         \textbf{Diotima: Compliance-Grade AI Infrastructure}\\[4pt]
         \textbf{for Formative Assessment in Education}}\\[10pt]

        % authors
        {\normalsize\scshape Noval Consulting}\\[4pt]

        % affiliation
        {\small\itshape
          AI Consulting and EU AI Act Advisory Practice,
          Ireland.\quad
          Contact: \texttt{info@novalconsulting.ai}}\\[10pt]

        % received line
        {\small (Working Paper, Draft v1, 2025)}\\[12pt]
      }

      \hrule height 0.4pt
      \vspace{10pt}

      % abstract
      \begin{abstract}
        \noindent
        The EU AI Act's classification of student assessment systems as high-risk under Annex~III
        creates both a compliance obligation and a strategic design question: what does an
        educational AI platform look like when it is built for regulatory alignment from first
        principles, rather than retrofitted to it? This paper describes Diotima, a
        teacher-centred formative assessment platform that treats Annex~III obligations not as
        constraints but as architectural foundations. By grounding all AI-generated content in
        approved curriculum materials, scoping high-risk components narrowly, enforcing human
        oversight as a workflow property rather than a policy statement, and embedding post-market
        monitoring as an operational discipline, Diotima demonstrates that high-risk AI in
        education can be simultaneously pedagogically sound, technically robust, and legally
        compliant. The paper argues that compliance-by-design is not a cost of participation in
        regulated education markets---it is the only credible path to adoption at scale.
      \end{abstract}

      \vspace{4pt}
      \noindent\textbf{Keywords}\quad
        EU AI Act\quad Annex~III\quad Formative assessment\quad
        Educational technology\quad Human oversight\quad Compliance-by-design

      \vspace{12pt}
      \hrule height 0.8pt
      \vspace{2pt}
      \hrule height 0.2pt
      \vspace{16pt}
    \end{@twocolumnfalse}
  ]%
}
\makeatother

% ─────────────────────────────────────────────────────────────────────────────

\begin{document}

\maketitle
\thispagestyle{fancy}

% ── INTRODUCTION ─────────────────────────────────────────────────────────────
\section{Introduction}

In proposing a new framework for AI governance in education, the EU AI Act draws
a conclusion that practitioners have long recognised implicitly: assessment is not a
neutral activity. \textsc{European Parliament} (2024, Art.~6) classifies AI systems
used to assess students in formal educational settings as high-risk under Annex~III,
Section~3(a). This classification reflects the reality that assessment shapes student
progression and opportunity, that biased or incorrect automated feedback can cause
material harm to learners, and that ungoverned automated evaluation erodes the professional
judgment on which good teaching depends.

Yet the proliferation of AI-powered assessment and feedback tools in the market suggests
that the full implications of this regulatory position have not been absorbed. Most
educational AI products are designed as if compliance were a downstream concern---a
documentation task to be completed once the product is built, rather than a constraint that
should shape every architectural decision from the outset.

The purpose of this paper is to describe a different approach. Diotima is a formative
assessment platform designed from first principles as a compliance-grade system---one that
treats Annex~III obligations not as a constraint on product development but as the
specification for it. The paper describes Diotima's architecture, its approach to human
oversight and explainability, its model governance framework, and its data governance
posture. It then examines the strategic implications of compliance-by-design as a
market positioning and institution-facing proposition.

The paper proceeds as follows. Section~2 establishes the regulatory context and the
specific obligations that Annex~III imposes. Section~3 describes Diotima's core design
proposition. Sections~4 through~9 examine the principal architectural and governance
dimensions in turn. Section~10 argues that compliance-by-design constitutes a durable
competitive differentiator in regulated education markets. Section~11 concludes.

% ── REGULATORY CONTEXT ───────────────────────────────────────────────────────
\section{AI in Education Is High-Risk by Default}

The EU AI Act's Annex~III classification of educational assessment AI as high-risk
reflects a straightforward chain of reasoning. Assessment shapes student progression
and opportunity. Biased or incorrect AI outputs can materially harm learners. And if
automated evaluation is allowed to operate without robust human governance, it
progressively displaces the professional judgment that should sit at the centre of any
assessment decision.

The high-risk designation applies even where AI is positioned as formative rather than
summative. Any system that analyses student responses, structures performance
interpretation, and generates feedback can meaningfully influence outcomes---regardless
of how it is labelled. The classification covers not just systems that assign official
grades, but any system whose outputs interact with the assessment process in a way that
could affect how a learner's performance is perceived or recorded (\textsc{European
Parliament}, 2024, Annex~III, s.~3(a)).

Diotima therefore adopts a conservative, regulator-aligned position from the outset: it
is treated as a high-risk system by design. That choice is not defensive---it is
strategic. Every architectural decision, every governance mechanism, every operational
workflow is oriented by that classification.

% ── CORE PROPOSITION ─────────────────────────────────────────────────────────
\section{Diotima's Core Design Proposition}

The most important thing to establish is what Diotima is \textit{not}. It is not an
AI grader. It does not autonomously assign marks, determine outcomes, or make binding
decisions about student performance. Diotima is a teacher-centred decision-support
system for formative assessment---its function is to help educators do their job
better, not to replace the professional judgment at the heart of good teaching.

In practice, that means Diotima supports teachers to: generate curriculum-aligned
questions grounded in approved source materials; apply consistent, transparent rubrics;
provide formative feedback at a scale that would otherwise be impossible; and support
student reflection and iterative improvement rather than delivering a single terminal
judgment.

The constraint that structures everything else: no AI output is final without human
confirmation. Teachers approve all assessment content. Teachers validate all evaluative
outputs. The AI proposes; the teacher decides. This positions Diotima as an assistive
infrastructure layer---one that raises the floor of feedback quality and frees teachers
to focus on the judgment calls that only a professional can make.

% ── SYSTEM OVERVIEW ──────────────────────────────────────────────────────────
\section{System Overview}

Diotima supports the complete formative assessment lifecycle from initial design through
to student reflection and audit. Every stage is engineered to support explainability,
traceability, and human oversight---not as retrospective additions, but as structural
properties of the workflow itself.

The lifecycle proceeds through seven stages: (1)~curriculum-grounded assessment design,
in which teachers select topics, subtopics, and Bloom's Taxonomy levels and the system
generates candidate content from approved source materials; (2)~teacher pre-approval of
every generated question, expected answer, and rubric before any student sees it;
(3)~student response and optional AI pre-check, allowing provisional analysis before
final submission; (4)~conditional rubric visibility, in which students receive
performance-band information after a genuine attempt, designed to support learning
without enabling gaming; (5)~teacher review, override, and finalisation, during which
all AI outputs remain provisional until confirmed by a teacher; (6)~student feedback
and reflection, in which final teacher-approved feedback is delivered and students are
prompted to engage with their learning; and (7)~comprehensive logging and auditability,
in which every action, inference, and decision is recorded to support traceability,
incident response, and post-market monitoring.

% ── ARCHITECTURE ─────────────────────────────────────────────────────────────
\section{Architecture Built for High-Risk Compliance}

\subsection{Grounded content generation.}

Hallucination is the most visible failure mode in educational AI.
Diotima addresses it architecturally rather than reactively. Every
AI-generated question, answer, and rubric is explicitly linked to curriculum
topics and subtopics, specified learning outcomes, Bloom's Taxonomy cognitive levels,
and approved, licensed source-of-truth materials.

This grounding serves three functions simultaneously. It prevents the model from
generating content that drifts from the curriculum. It ensures that every item can
be traced back to its source, giving teachers the visibility they need to make informed
approval decisions. And it provides a clear basis for explaining why an item exists and
what it is assessing---which is exactly what regulators and institutions require from
a high-risk system.

\subsection{High-risk components are narrowly scoped.}

A critical architectural decision is the deliberate separation of
high-risk components from supporting infrastructure. Only two components
are classified as high-risk: the Question, Answer and Rubric Generation Engine, and
the AI Inference and Feedback Engine. All supporting systems---the user interfaces,
workflow orchestration, data storage, and integration components---are explicitly
outside the high-risk perimeter.

This is not a trivial distinction. Narrowing the high-risk scope reduces risk
propagation, simplifies governance, and makes compliance obligations clearer and more
tractable for the teams responsible for maintaining them.

% ── HUMAN OVERSIGHT ──────────────────────────────────────────────────────────
\section{Human Oversight as a Structural Feature}

Article~14 of the EU AI Act requires meaningful human oversight for high-risk AI
systems. In most implementations, this ends up as a policy statement---a layer of
documentation applied over an architecture that was not designed with oversight in
mind. Diotima takes a different approach. Human oversight is not declared; it is
enforced by the workflow itself.

The structural properties are: teachers must approve every generated assessment item
before it reaches students; AI rubric placements are always provisional and are never
presented as final; teachers can accept, edit, override, or entirely replace any AI
output at any stage; and only teacher-confirmed results are stored as the authoritative
record. Students never receive AI-generated feedback that a teacher has not reviewed.
That is not a policy aspiration---it is a technical constraint embedded in the system.

The consequences of this design are significant. Professional autonomy is preserved.
The structural risk of over-reliance on AI outputs is prevented rather than mitigated
after the fact. And accountability stays where it belongs---with the human professional,
not the algorithm.

% ── EXPLAINABILITY ───────────────────────────────────────────────────────────
\section{Explainability for Teachers and Learners}

Explainability in Diotima is not a single mode---it is calibrated to role and context,
because what a teacher needs to understand is not the same as what a student needs to
understand, and the risks of over-transparency and under-transparency differ for each.

For teachers, the system provides full visibility into the rubric, the source text from
which assessment content was derived, and the AI's reasoning context. Every output can
be traced back to the curriculum materials that grounded it. The pattern of rejection
reasons and overrides constitutes an ongoing signal about model performance---a
practical feedback loop that informs model governance decisions.

For students, the explainability design is more deliberate. Rubric information is
revealed conditionally and progressively: students see performance bands after a
genuine attempt, not upfront. Full criteria are not disclosed in advance. This supports
transparency about how responses are being interpreted while preserving the pedagogical
integrity of the assessment, and it orients feedback toward improvement and next steps
rather than terminal judgment.

% ── MODEL GOVERNANCE ─────────────────────────────────────────────────────────
\section{Responsible Model Selection and Evaluation}

Diotima uses third-party pretrained foundation models rather than training its own.
That is a deliberate and defensible choice---but it does not mean model selection is
a technical decision. In Diotima, model choice is treated as a governance decision.

Models are evaluated and selected against a structured benchmark suite covering:
knowledge and reasoning (MMLU, GPQA); factuality and hallucination (LongFact);
instruction following (IFEval); bias and fairness (BBQ); and toxicity and safety
(Real Toxicity Prompts and custom safety evaluations).

No student data is used for model training at any stage. Model versions are logged
with every inference, creating the audit trail needed for incident investigation,
drift detection, and post-market monitoring. This approach enables evidence-based
model selection, ongoing comparison as new models emerge, and safe substitution of
model versions without architectural change---each of which matters for a system
expected to operate over a multi-year deployment horizon.

% ── DATA GOVERNANCE ──────────────────────────────────────────────────────────
\section{Data Governance and Privacy}

Data minimisation is a foundational principle in Diotima's design. The system
processes only what is necessary: student responses, assessment configuration, and
the metadata required to deliver formative feedback and support audit. There is no
collection of sensitive personal categories, no behavioural profiling, and no use of
student data to fine-tune or retrain the underlying models.

Retention limits are defined and aligned with institutional Data Protection Impact
Assessments. Educational content---curriculum documents, licensed textbooks, and
approved source materials---is managed under explicit publisher agreements covering
storage, processing, and derivative use within the scope of Diotima's purpose.

The effect is a data governance posture that protects every stakeholder simultaneously:
students, whose assessment data is not repurposed; teachers, whose professional actions
are logged for audit but not exposed inappropriately; publishers, whose content is
used under clear and documented licence; and institutions, whose DPIA obligations are
supported rather than complicated by the platform.

% ── POST-MARKET MONITORING ───────────────────────────────────────────────────
\section{Post-Market Monitoring and Continuous Assurance}

Compliance with the EU AI Act for a high-risk system does not end at deployment.
Post-market monitoring is a legal obligation---and in practice, it is where many AI
governance frameworks fail. Diotima embeds post-market monitoring from day one as an
operational discipline rather than an afterthought.

Models are regularly re-evaluated against the benchmark suite to detect drift and
performance degradation. Anomaly detection flags distribution shifts in AI outputs
before they reach users. Teachers have structured channels to report harmful or
incorrect outputs, with time-bound remediation workflows. The system is designed for
registration in the EU high-risk AI database, with version update processes in place.

The teacher approval and rejection workflow is itself a continuous monitoring
mechanism---every rejection reason is a signal about model performance in real
educational contexts that feeds back into the risk register and model governance
process. Compliance, in this design, is something the system generates evidence for
as a byproduct of normal operation.

% ── WHY COMPLIANCE IS THE MOAT ───────────────────────────────────────────────
\section{Why Compliance Is the Moat}

Most educational AI products treat regulation as friction to be minimised---a set of
boxes to check before getting on with the real work of building product. That framing
misreads what is happening in the market. Regulation is not friction. It is a filter.
And as the EU AI Act comes into force, it will progressively exclude products that were
not built to meet it.

Diotima treats compliance as infrastructure. By embedding Annex~III requirements into
the architecture rather than layering them on top of it, the platform achieves something
that post-hoc compliance cannot: genuine trustworthiness that institutions, teachers,
and regulators can verify rather than merely assert.

The downstream commercial logic is direct. Institutions can adopt AI-powered formative
feedback without regulatory anxiety, because the compliance case is built into what they
are acquiring. Teachers retain professional authority and the trust of their students,
because oversight is structural rather than nominal. Learners benefit from safer,
higher-quality feedback at a scale and consistency that manual processes cannot match.
And buyers gain access to markets---regulated school systems, post-secondary
institutions, professional certification bodies---that non-compliant products cannot
legally serve.

In an AI-regulated world, compliance-grade formative feedback is not optional. It is
the only scalable path forward.

% ── CONCLUSION ───────────────────────────────────────────────────────────────
\section{Conclusion}

The question facing every institution considering AI in education is not whether to
engage with the technology---that decision is already being made, whether or not formal
procurement processes have caught up with it. The question is whether to engage on
terms that are educationally sound, technically robust, legally defensible, and
professionally empowering for the teachers at the centre of it.

Diotima was designed to answer that question affirmatively. It demonstrates that
high-risk AI in education does not have to mean unsafe AI---that compliance with the
EU AI Act is not a ceiling on what the technology can do, but the foundation on which
trustworthy, scalable, pedagogically serious AI can be built.

This is not AI that replaces educators. It is AI that raises the floor of feedback
quality while keeping humans firmly in control. For institutions navigating the
regulatory moment in educational technology, that distinction is not incidental. It is
the entire point.

% ── REFERENCES ───────────────────────────────────────────────────────────────
\section*{References}
\addcontentsline{toc}{section}{References}

\begin{small}
\setlength{\parskip}{3pt}
\noindent
\textsc{European Parliament and Council of the European Union} (2024)
Regulation (EU) 2024/1689 of the European Parliament and of the Council of
13~June~2024 laying down harmonised rules on artificial intelligence (Artificial
Intelligence Act), \textit{Official Journal of the European Union}, L~2024/1689.

\medskip\noindent
\textsc{Noval Consulting} (2025)
Diotima AI Assessment Platform: Annex~III High-Risk AI Compliance Dossier, v1.1-draft.
Internal technical documentation, Noval Consulting, Dublin.

\medskip\noindent
\textsc{Krugman P.} (1991a)
Increasing returns and economic geography,
\textit{Journal of Political Economy} \textbf{99}, 483--499.

\medskip\noindent
\textsc{Bloom B. S., Engelhart M. D., Furst E. J., Hill W. H.} and \textsc{Krathwohl D. R.}
(1956) \textit{Taxonomy of Educational Objectives: The Classification of Educational
Goals. Handbook I: Cognitive Domain}. David McKay, New York.

\medskip\noindent
\textsc{Hendrycks D., Burns C., Basart S., Zou A., Mazeika M., Song D.}
and \textsc{Steinhardt J.} (2021) Aligning AI with shared human values,
\textit{arXiv preprint} arXiv:2008.02275.

\medskip\noindent
\textsc{Parrish A., Chen A., Nangia N., Padmakumar V., Phang J., Thompson J.,
Htut P. M.} and \textsc{Bowman S.} (2022) BBQ: A hand-built bias benchmark for
question answering, \textit{Findings of the Association for Computational
Linguistics: ACL 2022}, pp. 2982--2996.
\end{small}

\vspace{10pt}
\noindent\rule{\columnwidth}{0.4pt}\\
{\footnotesize\textit{Working paper derived from the Diotima Annex~III Technical
Dossier (v1.1-draft, 2025). Prepared by Noval Consulting. All correspondence to
\texttt{info@novalconsulting.ai}.}}

\end{document}
